% Edited conversion of source pp. 20--21.
% The three figure panels are extracted from the source PDF and decoded by LuaLaTeX.
\directlua{
  dofile("scripts/base64_asset.lua")
  decode_base64_file(
    "figures/07-triangular-plates/geometry.jpg.b64",
    "figures/07-triangular-plates/geometry.generated.jpg"
  )
  decode_base64_file(
    "figures/07-triangular-plates/mesh-folded.jpg.b64",
    "figures/07-triangular-plates/mesh-folded.generated.jpg"
  )
  decode_base64_file(
    "figures/07-triangular-plates/mesh-square.jpg.b64",
    "figures/07-triangular-plates/mesh-square.generated.jpg"
  )
}

\ReportHeading
  {Чисельний аналіз вільних коливань двох з’єднаних трикутних пластин}
  {М.~Ю. Борисенко\textsuperscript{1}, О.~В. Бойчук\textsuperscript{2}, Н.~П. Борейко\textsuperscript{1}, М.~Л. Тікіч\textsuperscript{1}}
  {\textsuperscript{1}Інститут механіки ім. С.~П. Тимошенка НАН України, Київ, Україна\\
   \textsuperscript{2}Миколаївський національний аграрний університет, Миколаїв, Україна}
  {}

З’єднані пластинчасті елементи різної геометричної форми, зокрема трикутні пластини, широко застосовують у тонкостінних конструкціях корпусів машин, надводних і підводних апаратів, літальних та космічних апаратів, а також інших інженерних споруд. Під час експлуатації такі конструкції зазнають статичних і динамічних навантажень різної природи, що зумовлює потребу в попередньому аналізі їх міцності, стійкості та динамічної поведінки.

Особливу увагу приділяють визначенню частот вільних коливань, оскільки резонансні режими можуть спричинити передчасне руйнування конструкції. У сучасній інженерній практиці динаміку пластинчастих систем аналізують за допомогою CAD/CAM/CAE-комплексів, що ґрунтуються на методі скінченних елементів. Одним із таких комплексів є Femap із розв’язувачем NX Nastran~\cite{plates:femap}. Цей підхід дає змогу з достатньою точністю визначати частоти та форми вільних коливань пластин складної геометрії~\cite{plates:hexagonal,plates:triangular}.

Мета роботи --- методом скінченних елементів визначити частоти та форми вільних коливань двох ізотропних тонких трикутних пластин, з’єднаних під двогранним кутом $\varphi$, за вільних країв та еквівалентної маси. За $\varphi=180^{\circ}$ дві трикутні пластини утворюють еквівалентну квадратну пластину.

\begin{figure}[htbp]
  \centering
  \begin{minipage}[b]{0.31\textwidth}
    \centering
    \includegraphics[width=\linewidth]{figures/07-triangular-plates/geometry.generated.jpg}\\[-0.2em]
    а)
  \end{minipage}\hfill
  \begin{minipage}[b]{0.31\textwidth}
    \centering
    \includegraphics[width=\linewidth]{figures/07-triangular-plates/mesh-folded.generated.jpg}\\[-0.2em]
    б)
  \end{minipage}\hfill
  \begin{minipage}[b]{0.31\textwidth}
    \centering
    \includegraphics[width=\linewidth]{figures/07-triangular-plates/mesh-square.generated.jpg}\\[-0.2em]
    в)
  \end{minipage}
  \caption{Геометрична модель двох з’єднаних трикутних пластин: а), б) --- загальний випадок; в) --- граничний випадок $\varphi=180^{\circ}$, що відповідає квадратній пластині.}
\end{figure}

\begin{thebibliography}{9}
\bibitem{plates:femap}
\emph{Dynamic Analysis User Guide: A Comprehensive User Guide for Simcenter Femap and Nastran Users}.
Portland, OR: Applied CAX, 2021. Available: \url{https://www.appliedcax.com/} (accessed Mar.~7, 2026).

\bibitem{plates:hexagonal}
A.~Grigorenko, M.~Borysenko, O.~Boychuk, and N.~Boreiko,
``Numerical analysis of free-vibration frequencies of a hexagonal plate,''
in \emph{Selected Problems of Solid Mechanics and Solving Methods},
Cham: Springer Nature Switzerland, 2024, pp.~201--220.

\bibitem{plates:triangular}
O.~Ya. Grigorenko, M.~Yu. Borisenko, O.~V. Boichuk, and L.~Yu. Vasil’eva,
``Free vibrations of triangular plates with a hole,''
\emph{International Applied Mechanics}, vol.~57, no.~5, pp.~534--542, 2021.
\end{thebibliography}
